% generated by roam.py handbook export -- не редактировать, править заметки
% 20260918131811-выпуклая_функция.org (body 061c43aff709)

\begin{definition}[{Выпуклая функция}]\label[definition]{note:1CA14F2C-637F-41DE-9E97-72075D19E97F}
  Пусть $E$~--- \hyperref[note:316B5BFA-8B9F-4267-9086-D6593AAB5F9B]{выпуклое множество}; $f: E \rightarrow \R$.
  $f$~--- выпуклая, если:
  \[\forall x,y \in E\ \forall \lambda \in [0;1]\ f(\lambda x + (1-\lambda)y) \le \lambda f(x) + (1-\lambda)f(y)\]

  Критерии и свойства:

  \begin{itemize}
    \item \hyperref[note:18BD8497-CEBD-4962-8144-474B489606C3]{Дифференциальный критерий выпуклости}
    \item \hyperref[note:86FCE16B-C82E-4416-851C-0E0AD0FDA382]{Критерий выпуклости в терминах вторых производных}
    \item \hyperref[note:1CD538BF-6F40-472A-A670-D38F093E4DF1]{Выпуклость множества подуровня выпуклой функции}
    \item \hyperref[note:035B8935-65A9-4BB4-BDA3-60A60FF278F0]{Наименьшее значение выпуклой функции достигается в точке с нулевым градиентом}
    \item \hyperref[note:6BC7D685-C429-45A3-B563-269840CA0664]{Преобразования, сохраняющие выпуклость}
  \end{itemize}
\end{definition}
